\section*{Bab \ref{ch:g11:minerals} --- Mineral}

\begin{solution}{exo:g11:minerals:1}
\omterm{def:g11:minerals:cleavage}{Belahan} adalah kecenderungan sebuah mineral untuk pecah di sepanjang bidang-bidang
datar, yang sejajar dengan arah tempat ikatan antara atom-atomnya lebih lemah;
\omterm{def:g11:minerals:cleavage}{pecahan} adalah patahnya mineral di sepanjang permukaan yang tidak teratur, pada
mineral yang tidak mempunyai bidang lemah. Kuarsa tidak mempunyai \omterm{def:g11:minerals:cleavage}{belahan}:
kuarsa memperlihatkan \omterm{def:g11:minerals:cleavage}{pecahan} konkoidal, dengan permukaan melengkung yang mulus
seperti kaca yang pecah. Mika mempunyai satu belahan yang sempurna: mika
terkelupas menjadi lembaran.
\end{solution}

\begin{solution}{exo:g11:minerals:2}
Sedikit unsur lain sudah cukup untuk mengubah warna kuarsa, yang dapat bening,
putih, hitam, merah muda, biru, atau ungu; olivin hampir selalu hijau. Yang lebih
dapat diandalkan: \omterm{def:g11:minerals:cleavage}{belahan} (atau \omterm{def:g11:minerals:cleavage}{pecahan}), kekerasan, \omterm{def:g11:minerals:lustre}{kilap}, \omterm{def:g11:minerals:habit}{habitus} kristal yang
tumbuh bebas, massa jenis.
\end{solution}

\begin{solution}{exo:g11:minerals:3}
(1) Olivin. (2) Biotit. (3) Ortoklas (sebuah feldspar: dua belahan, kekerasan
6). (4) Kuarsa (tanpa \omterm{def:g11:minerals:cleavage}{belahan}, kekerasan 7).
\end{solution}

\begin{solution}{exo:g11:minerals:4}
(1) Lima atom: satu silikon di pusat dan empat oksigen di sudut-sudut sebuah
tetrahedron. (2) $+4 + 4 \times (-2) = -4$: gugus itu adalah \ce{SiO4^4-}.
\end{solution}

\begin{solution}{exo:g11:minerals:5}
Olivin: tetrahedron terpisah. Piroksen: rantai tunggal. Amfibol: rantai ganda.
Mika: lembaran. Feldspar dan kuarsa: kerangka.
\end{solution}

\begin{solution}{exo:g11:minerals:6}
(1) $M(\ce{SiO2}) = 28.085 + 2 \times 15.999 = \qty{60.083}{g/mol}$;
oksigen: $31.998/60.083 \approx \qty{53.3}{\percent}$.
(2) $M(\ce{Mg2SiO4}) = 2 \times 24.305 + 28.085 + 4 \times 15.999 =
\qty{140.691}{g/mol}$; oksigen: $63.996/140.691 \approx \qty{45.5}{\percent}$.
(3) Kerak, dengan \qty{46}{\percent}, berorde sama: kerak hampir seluruhnya
terdiri atas mineral yang kaya akan oksigen, yaitu \omterm{def:g11:minerals:silicate}{silikat}.
\end{solution}

\begin{solution}{exo:g11:minerals:7}
Butir putih itu mempunyai dua belahan dan kekerasan antara 5.5 dan 6.5:
feldspar (6). Butir kelabu itu lebih keras daripada 6.5 dan memperlihatkan
\omterm{def:g11:minerals:cleavage}{pecahan} konkoidal yang melengkung tanpa \omterm{def:g11:minerals:cleavage}{belahan}: kuarsa (7).
\end{solution}

\begin{solution}{exo:g11:minerals:8}
(1) $4 - 2/2 = 3$ oksigen per silikon: \ce{SiO3}. (2) \ce{(Si,Al)2O6} memuat
dua atom silikon (atau aluminium) untuk enam oksigen, tiga per silikon: itu sama
dengan dua kali \ce{SiO3}.
\end{solution}

\begin{solution}{exo:g11:minerals:9}
$\rho = 0.55 \times 2.69 + 0.35 \times 3.4 + 0.10 \times 3.32 = 1.480 + 1.190
+ 0.332 \approx \qty{3.00}{g/cm^3}$. Kubus itu bervolume
\qty{1000}{cm^3}, jadi massanya sekitar \qty{3.00}{kg}.
\end{solution}

\begin{solution}{exo:g11:minerals:10}
(1) $0.75 \times 3.32 + 0.25 \times 3.4 = 2.49 + 0.85 = \qty{3.34}{g/cm^3}$.
(2) $(3.34 - 3.00)/3.00 \approx 0.11$: sekitar \qty{11}{\percent} lebih rapat.
(3) Basalt dan \omterm{def:g11:ocean-ridges:gabbro}{gabro} \omterm{def:g11:ocean-ridges:oceanic-lithosphere}{kerak samudra} lebih ringan daripada \omterm{def:g11:ocean-ridges:peridotite}{peridotit} mantel:
keduanya tetap berada di atasnya.
\end{solution}

\begin{solution}{exo:g11:minerals:11}
(1) Plagioklas $168/400 = 0.42$, piroksen $0.33$, olivin $0.15$, lubang
$0.10$. (2) Proposisi itu berlaku untuk batuan tanpa pori; di sini sepersepuluh
volumenya kosong. (3) Pada bagian yang padat, bagian-bagiannya adalah $0.42/0.90
\approx 0.467$, $0.367$, dan $0.167$:
$\rho_{\text{padat}} \approx 0.467 \times 2.69 + 0.367 \times 3.4 + 0.167
\times 3.32 \approx \qty{3.06}{g/cm^3}$. Dengan lubang-lubangnya, yang menambah
volume tanpa menambah massa, batuan itu bermassa jenis $0.42 \times 2.69 + 0.33 \times 3.4 +
0.15 \times 3.32 \approx \qty{2.75}{g/cm^3}$, yaitu $0.9 \times 3.06$.
\end{solution}

\begin{solution}{exo:g11:minerals:12}
(1) Oksigen $464/16.0 = 29.0$~mol; silikon $283/28.1 \approx 10.1$; aluminium
$84/27.0 \approx 3.1$; besi $52/55.8 \approx 0.93$; kalsium $46/40.1 \approx
1.15$; magnesium $28/24.3 \approx 1.15$; natrium $23/23.0 = 1.0$; kalium
$15/39.1 \approx 0.38$. (2) Seluruhnya sekitar $46.8$~mol atom, dengan
$29.0/46.8 \approx 0.62$ di antaranya oksigen: hampir dua dari tiga atom adalah
oksigen. (3) Ion-ion oksigen yang besar membentuk sebagian besar atom, dan lebih
besar lagi bagiannya dalam volume: pada skala atom, kerak adalah susunan rapat
ion-ion oksigen, dengan ion-ion silikon dan logam yang kecil di ruang-ruang di
antaranya.
\end{solution}

\begin{solution}{exo:g11:minerals:13}
(1) Banyaknya rata-rata sudut yang dipakai bersama adalah $(2 + 3)/2 = 2.5$, jadi
ada $4 - 2.5/2 = 2.75$ oksigen per silikon. (2) Empat silikon untuk sebelas
oksigen: \ce{Si4O11}. (3) \ce{Si7AlO22} memuat delapan tetrahedron (tujuh
silikon, satu aluminium) untuk 22 oksigen, yaitu $22/8 = 2.75$ per tetrahedron,
atau dua kali \ce{Si4O11}. Kedua gugus \ce{OH} bukan bagian dari
tetrahedron-tetrahedron itu.
\end{solution}

\begin{solution}{exo:g11:minerals:14}
$0.20 \times 2.62 + 0.10 \times 2.56 + 0.50 \times 2.65 + 0.10 \times 3.09 +
0.10 \times 3.23 = 0.524 + 0.256 + 1.325 + 0.309 + 0.323 \approx
\qty{2.74}{g/cm^3}$, di atas \qty{2.7}{g/cm^3}; granit pada contoh itu, yang
lebih kaya akan kuarsa dan feldspar kalium, sekitar \qty{2.59}{g/cm^3}, di bawah
\qty{2.6}{g/cm^3}. Batuan di bagian barat lebih rapat karena memuat lebih banyak
plagioklas serta lebih banyak biotit dan hornblende, yang kaya akan besi dan
magnesium; granit yang tinggi kadar silikanya memuat lebih banyak kuarsa dan
feldspar kalium, mineral pembentuk batuan yang paling ringan.
\end{solution}

\begin{solution}{exo:g11:minerals:15}
(1) Kerak atas: $66.6 \times 0.467 \approx \qty{31.1}{\percent}$ silikon;
seluruh kerak: $60.6 \times 0.467 \approx \qty{28.3}{\percent}$. (2) Kerak atas,
yang lebih kaya akan silika, seharusnya mengandung lebih banyak kuarsa dan
feldspar, dan lebih sedikit mineral yang kaya akan besi dan magnesium; kerak yang
lebih dalam, sebaliknya. (3) Kuarsa dan feldspar adalah mineral pembentuk batuan
yang paling ringan: kerak atas seharusnya lebih ringan daripada kerak yang lebih
dalam.
\end{solution}

\begin{solution}{pb:g11:minerals:1}
\textbf{1.} Kuarsa: tanpa \omterm{def:g11:minerals:cleavage}{belahan}, \omterm{def:g11:minerals:cleavage}{pecahan} konkoidal, \omterm{def:g11:minerals:lustre}{kilap} kaca, lebih keras
daripada paku (7).

\textbf{2.} Ortoklas, sebuah feldspar kalium: dua belahan, kekerasan 6 pada
skala Mohs.

\textbf{3.} Plagioklas (di sini oligoklas), feldspar yang mengandung natrium dan
kalsium.

\textbf{4.} D adalah biotit, E muskovit: dua jenis mika. Tetrahedronnya membentuk
lembaran, dan ikatan antara lembaran-lembaran itu lemah: keduanya terbelah di
sepanjang \omterm{def:g11:minerals:cleavage}{belahan} tunggal yang sempurna ini.

\textbf{5.} Kuarsa dan kedua feldspar (A, B, C) adalah kerangka; biotit
(D) mengandung besi dan magnesium.

\textbf{6.} Kuarsa $150/500 = 0.30$; ortoklas $0.35$; oligoklas $0.25$;
biotit $0.06$; muskovit $0.04$.

\textbf{7.} Bila butir-butirnya tersebar tanpa keteraturan, bagian titik pada
sebuah mineral menaksir bagian luas yang ditutupinya pada permukaan itu, dan
bagian itu sama dengan bagian volume yang diisinya.

\textbf{8.} $0.35 + 0.25 = 0.60$: \qty{60}{\percent} feldspar.

\textbf{9.} Sebesar $1/500 = 0.002$, yaitu \qty{0.2}{\percent}.

\textbf{10.} Ya: feldspar, kuarsa, dan mika.

\textbf{11.} $\rho = \phi_1 \rho_1 + \phi_2 \rho_2 + \dots + \phi_k \rho_k$,
dengan $\phi_1 + \dots + \phi_k = 1$.

\textbf{12.} Kuarsa $0.30 \times 2.62 = 0.786$; ortoklas $0.35 \times 2.56 =
0.896$; oligoklas $0.25 \times 2.65 = 0.6625$; biotit $0.06 \times 3.09
\approx 0.185$; muskovit $0.04 \times 2.82 \approx 0.113$ (semuanya dalam
\unit{g/cm^3}).

\textbf{13.} $0.786 + 0.896 + 0.6625 + 0.185 + 0.113 \approx
\qty{2.64}{g/cm^3}$.

\textbf{14.} Ortoklas, dengan $0.896$, karena mineral itu yang paling melimpah.
Biotit dan muskovit, dengan \qty{3.09}{} dan \qty{2.82}{g/cm^3}, lebih rapat
daripada mineral-mineral lain dan menaikkan massa jenis batuan di atas massa jenis
kuarsa dan feldspar saja.

\textbf{15.} \qty{2640}{kg/m^3}: \qty{1}{m^3} granit itu bermassa
sekitar \qty{2.64}{t}.

\textbf{16.} $100 \times 2.643 \approx \qty{264.3}{g}$.

\textbf{17.} Blok itu mendesak \qty{100}{cm^3} air, yang bermassa
\qty{100}{g}: timbangan menunjukkan $264.3 - 100 = \qty{164.3}{g}$.

\textbf{18.} $264.3/(264.3 - 164.3) = 264.3/100 \approx \qty{2.64}{g/cm^3}$.

\textbf{19.} Di antara keduanya: lebih ringan daripada rata-rata
\qty{2.69}{g/cm^3}, di atas \qty{2.6}{g/cm^3} granit yang paling kaya akan silika
dan kalium. Lebih banyak plagioklas dan lebih banyak biotit atau hornblende, serta
lebih sedikit kuarsa dan feldspar kalium, akan membuat sebuah granit lebih rapat.

\textbf{20.} Massa jenis granit itu sekitar
\textbf{\qty{2.64}{g/cm^3}}.
\end{solution}
